\section{A fixed surface tuple on a nonempty outer divisor collar}
\label{sec:r13-fixed-surface-tuple-collar}

Keep the original finite extension \(K/\mathbf Q_p\), its norm,
and the divisor tube of DT and AF.  The central model has
dimension three, the supplied divisor is \(D\), its actual
surface factor is \((Y_D,v)\), and its normal parameter is
\(Z\).  The surface point has the original discrete value
group.  Write \(q_D(r)\) for the actual relative Gauss point
with \(|Z|=r\).  All ordinary domains below are domains in
these original analytic spaces.

\subsection{One finite set of controls for every outer annulus}

The proof of AF1 uses a finite list of zero-slice functions
on an ordinary surface neighborhood \(W_0\ni v\).  Its
invertible coefficient and localizer functions \(u_i\) have
\(|u_i(v)|=1\); its other coefficient functions \(a_j\) have
\(|a_j(v)|\leq1\).  Inverse functions needed in the estimates
are included among the \(u_i^{\pm1}\).  This list is fixed by
the given finite central conditions and standard etale chart.
For every compact interval \(I\Subset(r_*,1)\), the argument
chooses a number \(c_I>1\), sufficiently close to one, and
uses only the surface restrictions
\[
       c_I^{-1}<|u_i|<c_I,\qquad |a_j|<c_I.
\tag{UA1}
\]
The constants \(M\), the finitely many divisor orders, and
the preliminary threshold \(r_*<1\) are independent of
\(I\).  Bounds such as \(r_+c_I^{M+2}<1\) determine how
close \(c_I\) is to one; they do not change the list of
functions.

Choose a finite tuple of surface units \(C\) as follows.
Retain every \(u_i\).  If \(|a_j(v)|=1\), retain \(a_j\),
after restricting its unit domain.  If \(|a_j(v)|<1\),
retain \(1+a_j\).  Every coordinate of \(C\) has norm one
at \(v\).  Shrink \(W_0\) once so all these are actual
units.  Put \(t_C=-\log|C|\), coordinatewise.

\begin{lemma}[Fixed controls and a fixed protected central fiber]
\label{lem:r13-fixed-controls-protected-fiber}
Let \(F\) be any fixed finite surface tuple retaining \(C\),
and apply GP once on a prescribed ordinary neighborhood of
\(v\) in \(W_0\), retaining \(F\).  Denote the protected
tuple by \(G\), its central value by \(b=t_G(v)\), and its
compact entire central fiber by \(E=t_G^{-1}(b)\).  Then:
\begin{enumerate}
\item Every point of \(E\) satisfies \(|u_i|=1\) and
\(|a_j|\leq1\).
\item For every \(c>1\), one may narrow the target bands
around \(b\), without changing any unit, so their complete
ordinary inverse image satisfies
\(c^{-1}<|u_i|<c\) and \(|a_j|<c\).  The closed narrower
bands still have compact inverse image in the original
interior.  The entire central fiber and all its source
germs remain unchanged.
\item For every \(I\Subset(r_*,1)\), the actual AF product
isomorphism exists on a product containing the relative
Gauss family over all of \(E\), with the narrowed surface
bands of the second assertion.  It satisfies the original
finite central conditions at every analytic product point.
\end{enumerate}
\end{lemma}

\begin{proof}
All coordinates of \(C\) have their norm-one values on
\(E\).  In the case \(|a_j(v)|<1\), the ultrametric
inequality gives
\[
                |a_j|=|(1+a_j)-1|\leq1
                       \quad\hbox{on }E.
\tag{UA2}
\]
The other assertions in the first item follow directly
from the retained coordinates.  To prove the second item,
choose \(1<c'<c\) and restrict the \(C\)-coordinate bands
to \((c'^{-1},c')\).  For a shifted coordinate, its upper
bound gives \(|a_j|\leq\max(|1+a_j|,1)<c'\).  Use still
smaller bands for the closed target closure if necessary.
All these bands contain \(b\), so they retain exactly the
same entire central fiber.  GP supplies compactness and
the interior property before narrowing, and these properties
persist.  Source germs depend on the tuple and an ordinary
neighborhood of their points, not on the widths of bands
containing their central value.

For the last assertion use the AF estimates with the same
finite functions, but on a neighborhood of \(E\) satisfying
UA1.  They are inequalities for zero-slice analytic functions
and hold at every analytic product point.  For an open
interval \(J\) with \(I\Subset J\Subset(r_*,1)\), choose
\(c_I\) so that all strict Newton and divisor-domination
inequalities AF4--AF8 hold.  Cover the compact surface set
by finitely many affinoid neighborhoods contained in the
allowed narrower bands.  Newton iteration on the products
gives the selected root; uniqueness in its Hensel ball
glues the roots and the inverse zero-slice maps.  Taking
ordinary interiors gives the required product.  The same
finite basic central inequalities AF2 place its entire
image in the prescribed original ordinary domain.  This
proof uses the fixed bounds on \(E\); it does not require
that an arbitrary surface neighborhood contain \(E\).
\end{proof}

\subsection{The interval-independent finite preparation}

Fix the initial finite surface coefficient tuple of AF9,
of full normal rank two, and adjoin \(C\).  Keep the fixed
finite closed-cell PS hypotheses of NSA: smooth extensions
near every closed coefficient cell, compatible prescribed
face restrictions, and bounded first derivatives on each
smaller compact truncation.  Bounds are taken only after
an annular interval and truncation have been chosen.

\begin{proposition}[Finite surface preparation before the radial interval]
\label{prop:r13-surface-preparation-before-interval}
In this fixed AF setting, the unit tuple in the annular
NSA construction can be chosen before the compact radial
interval.  More precisely, there is one finite surface
tuple \(H\), consisting of units of norm one at \(v\)
after original-constant normalization, and one fixed
finite sequence of surface geometric preparations, with
\(H\) the last protected tuple on its fixed surface
unit domain.  Write \(E_H=t_H^{-1}(t_H(v))\) for its
entire compact central fiber.  These data have
the following property.  For every
\(I\Subset(r_*,1)\), the construction yields a continuous
finite-PS ambient covector \(\Lambda_I\) in the coordinates
\((Z,H)\) on an ordinary neighborhood of the entire
product over \(I\) and \(E_H\), with
\[
                    t_{(Z,H)}^*e(\Lambda_I)=T,
                       \qquad e=\operatorname{sym}d,
\tag{UA3}
\]
whenever the represented tensor satisfies the original
current equation on the AF product.  This is the complete
ordinary image identity, with all its lower faces.  The
tuple \(H\) and the geometric choices do not depend on
\(I\); the coefficient operators, ordinary truncations
and bounds may depend on \(I\).
\end{proposition}

\begin{proof}
We specify the choices in the order in which their data
become available.  Every protection performed in this
proof retains \(C\).  Thus Lemma~\ref{lem:r13-fixed-controls-protected-fiber}
permits it to be made once on the surface and used in
every later annular interval by changing only target bands.
Every ordinary localization used for residue identities,
unit domains or full-rank fiber isolation is a fixed
surface localization.  It can be imposed before that
protection.  All analytic domains are common domains of
finitely many such requirements.

\emph{Entry and the full trace space.}
Perform the first CE protection for the fixed old surface
tuple.  Its complete finite central source fiber and every
projected satellite ray are now fixed.  Insert those rays
once to obtain CE's rational normal fan \(\Sigma_0\).
In particular this subdivision is not recomputed when
the tangent interval changes.  The finite-dimensional
space of continuous homogeneous functions linear on
\(\Sigma_0\) is defined by rational face equations.
Choose CG's rational ambient PL/DC extension basis for
this entire space, rather than for the span of the
coefficient functions on one particular interval.

SC applies to that fixed rational family and constructs
its actual cap units, numerical matrix, rational kernel
basis and RB unit representatives before any coefficient
vector is chosen.  All these functions therefore can be
chosen now.  Adjoin them and make the next fixed surface
localization and protection required for SS.  These steps
fix the old tuple \(\tau\) for NSA.  On each later
interval CE and SP construct the conewise primitives;
CG's affine projector is chosen on that interval.  The
resulting central trace lies in the already prepared
full space.  SS and CU use the same fixed cap and unit
families to normalize it.  Their arbitrary compact
tangent tests and the complete protected source sums
are legitimate by the preceding lemma.  Changing the
affine projector changes coefficients, not the PL
basis, caps or units.

\emph{Core, source separation and transverse bases.}
Apply SA to the frozen \(\tau\), then BG, SN and the
source-exhaustion theorem.  Their smooth proper surface,
finite boundary graph, actual node coordinates and
selected-factor localizers depend only on that tuple.
Choose the first FA family on its complete full source
star, with the finite algebraization coordinates retained.
Make the corresponding fixed surface localization and
protection.  Choose the full finite marked sets on all
tracked divisor curves, and choose RT's Picard basis for
the entire weighted zero-sum space on each such set.
SL lifts these functions simultaneously, with residue
one on the other tracked divisors.  Include the finite
SA coordinates required for these functions.

No value of a primitive or a trace function is used in
these choices.  The weighted spanning tree can likewise
be fixed now.  On a later interval RCF computes the
sourcewise flux coefficients, and the tree chooses the
rotation coefficients.  RT is its fixed real linear
map on the resulting entire trace family.  Thus none
of these differential steps adds a unit.  The actual
new-mark and primitive-index checks in NSA are checks
of these fixed finite coordinate maps.  They apply on
every interval with the same labels and weights.

The intermediate protection coordinates cause no new
coefficient-dependent algebraic choice: their norms
are constant on an ordinary surface germ at \(v\).
They can be retained in the fixed pre-Picard tuple,
as NSA does.  Their selected main-source residues are
constants after normalization.  Their additional
protected source germs and projected rays have already
been fixed by the corresponding protection; the NSA
satellite argument retains them in each source sum.

\emph{Final preparation.}
Freeze the tuple containing the preceding coefficients,
caps, kernel units, nodes, first separators, Picard units
and algebraization coordinates, and retain every earlier
protection coordinate before the last protection.
Choose the final FA/CS
separating family and any finite SA completion.  Localize
the complete old full-rank central fiber once.  Perform
the last surface protection, retaining \(C\) again.
Call this last protected tuple \(H\).  All earlier
coordinates have already been retained; there is no
additional coordinate choice after this protection.

For a specified interval, narrow the finitely many
protected target bands using UA1 and the lemma, take
a common smaller normal truncation, and run NSA's
coefficient operations on a real tangent interval
containing \(-\log I\).  The complete image proof of
NSA checks every actual incidence, including the
additional protected sources.  The narrowed target
bands contain \(t_H(v)\), hence their full inverse
image contains the entire same compact fiber \(E_H\).
The fixed-control lemma supplies an AF product over
that whole fiber for the chosen interval.  Its final NL extension
uses only the fixed finite coordinate spaces; the
closed-cell coefficient subdivision and its analytic
extensions may be restricted to the chosen interval.
NL does not select a new analytic unit.  This proves
UA3.  Intersecting the real coefficient cells with a
new tangent interval or with a fixed geometric cone
changes a real subdivision only; no unit is selected
from those coefficient walls.  There are finitely
many geometric stages, and
all have been fixed before the interval.  The full
fixed-tuple assertion follows.
\end{proof}

\subsection{Central lift first, then every sufficiently outer interval}

\begin{theorem}[A nonempty central-compatible punctured collar]
\label{thm:r13-central-compatible-punctured-collar}
Let a fixed central finite coefficient presentation on
the original threefold have the AF closed-cell PS
pullback, and assume \(D[T]=0\) on its original ordinary
domain against all original compact smooth complementary
CLD tests.  Then there are one finite actual central
unit tuple \(\widehat H\), retaining the original
coefficient tuple and \(Z\), and a number \(r_0<1\),
such that the following holds for every compact
\(I\Subset(r_0,1)\).

There is an ordinary neighborhood \(U_I\) of the
entire actual interval \(q_D(I)\), and a continuous
finite-PS ambient covector \(\widehat\Lambda_I\)
in the fixed \(\widehat H\)-coordinates, with smooth
closed-cell extensions and bounded first derivatives
on a smaller fixed truncation, whose strain pulls
back to \(T\) on the complete ordinary image and
every lower face.  Consequently there are smooth
full-exact forms
\[
       a_{I,\eta}=d'\beta_{I,\eta}+d''\zeta_{I,\eta},
              \qquad \epsilon_{I,\eta}\longrightarrow0,
\tag{UA4}
\]
and one fixed positive control \(Q_I=D[q_I]\) on
that domain satisfying the original NSA4 order and
all-background inequalities.  For each \(I\), the
tuple, covector, control and approximation precede
every later allowed symmetric middle-positive
separately \(d'\)- and \(d''\)-closed background,
every original compact complementary smooth test,
and every actual refinement, including rank loss.
The central tuple \(\widehat H\) is independent of
\(I\).
\end{theorem}

\begin{proof}
Prepare the fixed surface tuple \(H\) of the preceding
proposition.  Normalize its central surface values
by constants in the original \(K\).  For each of its
finitely many residue functions in \(k(D)^\times\),
choose a unit of the residue local ring at \(D\)
mapping to it, and lift that function to an actual
unit at the original threefold point using the
finite same-field unit lifting theorem.  These
central units have divisor order zero at \(D\).
Use DT8 for their exact norm and normalized
initial-form comparison with the prescribed
surface units; the finite ordinary comparison
of two surface lifts with the same residue is
included as well.  The construction uses the
original field and its original norm throughout.

Apply AF1 to this now fixed enlarged central list.
It supplies a threshold \(r_0<1\), enlarged to
include the initial \(r_*\), and for every
\(I\Subset(r_0,1)\) an actual product domain
\(\mathcal A_J\times W_I^{\rm lift}\), with
\(I\Subset J\) and \(v\in W_I^{\rm lift}\),
on which all surface coordinates in \(H\) have
the exact norms of their central lifts.  The
original coefficient coordinates retain their
fixed integral affine graph formulas involving
\(Z\).  Every radial value is included.

Run the preceding proposition on the same
interval, using the original finite central
conditions.  It gives \(\Lambda_I\) on a
product neighborhood
\(\mathcal A_{J'}\times W_I^{\rm strain}\)
of the relative Gauss interval.  Shrink the
two open radial intervals so both still
contain \(I\).  The ordinary surface
intersection
\[
       W_I^{\rm lift}\cap W_I^{\rm strain}
                          \ni v
\tag{UA5}
\]
is nonempty.  Both products use the same
selected zero-slice comparison, by Newton
uniqueness.  Hence they intersect in an
actual product containing \(q_D(I)\).
This is the required nonvacuity: the interval
was chosen after the final central-lift
threshold, and no new coefficient unit is
selected afterward.  Extra finite controls
used to describe the lift domain impose
only an ordinary restriction; they are not
inputs to another scalar trace construction.

On this intersection the exact coordinate
identities transport \(\Lambda_I\) to
\(\widehat\Lambda_I\).  Pullback through
their fixed integral affine graph gives
the full strain identity on each actual
face, with the prescribed lower-face
values.  Every later actual refinement
is a linear coordinate pullback after
finite common subdivision, so rank loss
does not change this identity.  Restrict
to a smaller ordinary protected truncation
and apply the retained ambient strain
smoothing theorem.  It gives UA4 and
exactly the control and all-background
inequalities of NSA4, with their original
tests and positivity hypotheses.
\end{proof}

\paragraph{Quantifiers.}
The theorem supplies one central tuple and a
primitive for every compact interval in an
actual nonempty punctured collar.  The
primitives, domains and finite constants can
depend on the interval.  Compatibility
between those primitives at the original
apex is a further coefficient problem.
Here the complete image is the image of
the displayed ordinary domain \(U_I\).
Other actual sources in a larger central
domain with the same completed moment
value over the old ray are not included
unless they belong to \(U_I\).  Exhausting
their complete incident family is a
separate source-localization question.

